\section{Related work}\label{sec:related}

\paragraph{Bivariate direction rules.} The identifiability results for linear non-Gaussian models \citep{shimizu2006linear}, additive-noise models \citep{hoyer2009nonlinear,peters2014identifiability,zhang2009identifiability} and post-nonlinear models underlie most pair rules. IGCI \citep{daniusis2010inferring} rests on the independence between the cause distribution and the mechanism; RECI \citep{blobaum2018cause} compares regression errors after scaling to the unit interval; pairwise LiNGAM statistics \citep{hyvarinen2013pairwise} compare likelihoods of non-Gaussian linear models. LOCI \citep{immer2023loci} fits location-scale noise models with neural networks and an independence test. SLOPE \citep{marx2017slope} uses minimum description length with a regression-based model class. Learned approaches train a classifier on simulated pairs \citep{lopezpaz2015towards}; AVICI \citep{lorch2022amortized} amortises structure inference over simulated data and is the learned graph method we run.

\paragraph{Evaluation of causal discovery.} Real data with known cause--effect structure are scarce. The T\"ubingen collection \citep{mooij2016distinguishing} is the standard for pairs; the Causal Chambers \citep{gamella2025chamber} provide physical devices with a documented graph and interventional recordings; the Sachs network \citep{sachs2005causal} is a consensus rather than a ground truth. For graphs on simulated data, \citet{reisach2021beware} showed that the variance structure of simulated additive-noise models can be read by a trivial rule, which then rivals score-based methods; our marginal-only classifier is a pair-level analogue of that observation on real data. A same-hands re-evaluation on the T\"ubingen pairs, which also introduced the compression baseline we use, is \citet{stienstra2026oneruler}; we reproduce its main figure in Section~\ref{sec:audit}.

\paragraph{Recent graph methods.} BF-BIC with BOSS \citep{ramsey2025bfbic} scores nonlinear recursive models with truncated basis functions; SNOE \citep{huang2025snoe} orients edges with pairwise additive-noise likelihoods; topological residual asymmetry \citep{elbouchattaoui2026tra} compares residual clouds. We implement versions of all three from their descriptions where public code was not available, and say where our constants are our own (Section~\ref{sec:data}).

\paragraph{Reporting and statistical resolution.} The concern that benchmark comparisons are underpowered and that reporting choices are free parameters is general in machine learning and statistics; we do not review it. What is specific here is that the quantities are measured on causal-discovery benchmarks, with exact paired tests on stored decisions and an enumerated multiverse of choices that have each been used for pairs.
